\chapter{MRI}

\section*{Definition and Overview}
Magnetic Resonance Imaging (MRI) is a non-invasive, high-resolution diagnostic imaging modality that uses strong magnetic fields, radiofrequency pulses, and advanced computer algorithms to generate detailed cross-sectional anatomical images of the human body. Unlike plain radiography or computed tomography (CT), MRI does not utilise ionising radiation. It provides exceptional soft-tissue contrast, making it particularly valuable for evaluating the brain, spinal cord, musculoskeletal joints, cardiovascular system, and pelvic organs. MRI operates on the principles of nuclear magnetic resonance, specifically targeting hydrogen protons in water molecules within bodily tissues. Despite its high safety profile, the strong magnetic field necessitates rigorous pre-screening for metal implants, pacemakers, or foreign bodies.

\section*{Epidemiology}
Millions of MRI scans are performed globally each year, representing a cornerstone of modern diagnostic medicine. In Australia, over one million Medicare-subsidised MRI scans are performed annually, with usage rising steadily due to expanded indications and ageing populations. The most common anatomical regions scanned are the brain and spine (accounting for approximately 50\% of scans), followed by musculoskeletal imaging (knees, shoulders) and oncology staging. While MRI is safe for patients of all ages, including pregnant individuals (typically after the first trimester), access is regulated by high equipment costs, specialist referral requirements, and safety contraindications.

\section*{Aetiology and Pathophysiology}
While not a disease, the clinical indications for MRI and its underlying physical principles are rooted in human tissue composition and pathophysiology:
\begin{itemize}
    \item \textbf{Hydrogen proton alignment:} When a patient enters the scanner, the strong static magnetic field (measured in Tesla, typically 1.5T or 3T) aligns the spin of hydrogen protons (which act like tiny magnets) parallel or antiparallel to the field.
    \item \textbf{Radiofrequency excitation:} A temporary radiofrequency (RF) pulse is applied, tilting the protons out of alignment. Once the RF pulse is turned off, the protons return to their original alignment (relaxation phase).
    \item \textbf{T1 and T2 relaxation signals:} As protons relax, they emit RF signals detected by receiver coils. T1-weighted images show fat as bright and fluid as dark (ideal for anatomy). T2-weighted images show fluid as bright (ideal for identifying pathological oedema or inflammation).
    \item \textbf{Pathological soft-tissue changes:} MRI detects minute variations in water content and proton mobility. Pathological processes such as demyelination (multiple sclerosis), ischemia (acute stroke), neoplasia (tumours), or trauma (ligament tears) alter these properties, generating distinct signals.
    \item \textbf{Gadolinium contrast enhancement:} Intravenous gadolinium-based contrast agents alter local magnetic properties, highlighting areas of increased vascularity or blood-brain barrier disruption.
\end{itemize}

\section*{Clinical Features}
Features of the patient experience and clinical indicators that point to the need for an MRI scan include:
\begin{itemize}
    \item \textbf{Neurological symptoms:} Persistent headaches, progressive sensory or motor deficits, cognitive decline, seizures, or sudden visual field deficits.
    \item \textbf{Chronic joint pain or instability:} Suspected ligament tears (e.g. ACL), meniscus tears, or tendonitis refractory to conservative management.
    \item \textbf{Radicular pain:} Radiating pain, numbness, or weakness in the arms or legs, suggestive of spinal disc herniation or spinal canal stenosis.
    \item \textbf{Oncological staging features:} Suspected primary tumours or metastatic disease requiring high-definition soft-tissue characterisation.
    \item \textbf{Pelvic mass symptoms:} Chronic pelvic pain or abnormal bleeding where ultrasound is inconclusive (evaluating uterine fibroids, adenomyosis, or ovarian pathology).
    \item \textbf{Claustrophobia during screening:} The patient experiencing intense anxiety or panic due to the narrow bore and loud noise of the scanner, requiring sedation or open-bore alternatives.
\end{itemize}

\section*{Differential Diagnosis}
Clinicians must choose the appropriate imaging modality depending on the diagnostic requirements:
\begin{itemize}
    \item \textbf{MRI vs. CT scan:} CT is faster, uses radiation, and is superior for acute intracranial haemorrhage, bone fractures, and pulmonary embolisms. MRI is slower (20--60 minutes) but offers far superior detail for soft tissues, brain stem lesions, and spinal cord pathology.
    \item \textbf{MRI vs. Ultrasound:} Ultrasound is dynamic, cheaper, and portable, making it ideal for superficial tendons, obstetric imaging, and vascular flow, whereas MRI provides comprehensive structural detail deep within joints and cavities.
    \item \textbf{MRI vs. X-ray:} X-rays are the primary screen for bone pathology (fractures, osteoarthritis), whereas MRI is reserved to evaluate soft-tissue components like cartilage, ligaments, or marrow infiltration.
    \item \textbf{Gadolinium reaction vs. Iodinated contrast reaction:} Gadolinium contrast carries a much lower risk of severe allergic reactions compared to the iodinated contrast used in CT scans.
    \item \textbf{Nephrogenic Systemic Fibrosis (NSF) vs. Other skin sclerosing disorders:} A rare, severe fibrotic complication of gadolinium in patients with severe renal failure, which must be distinguished from systemic sclerosis.
\end{itemize}

\section*{When to Seek Medical Care}
Patients scheduled for an MRI must report any metal implants, medical devices, or severe claustrophobia to their healthcare provider.

\textbf{Call 000 (or local emergency services) for:}
\begin{itemize}
    \item \textbf{Acute stroke symptoms:} Sudden weakness on one side of the body, facial drooping, or difficulty speaking, which requires urgent emergency stroke protocol imaging (typically CT or rapid MRI).
    \item \textbf{Acute spinal cord compression:} Sudden loss of bladder or bowel control, accompanied by numbness in the groin or buttocks (saddle anaesthesia) and leg weakness (cauda equina syndrome).
    \item \textbf{Anaphylaxis to contrast:} Sudden breathing difficulty, swelling of the throat, or widespread hives immediately after receiving intravenous gadolinium.
\end{itemize}
Other reasons to contact the imaging clinic before your scan include having a pacemaker, cochlear implant, programmable shunt, metal fragments in the eyes, or severe kidney disease.

\section*{Diagnosis and Investigations}
Safety screening and preparation are vital components of the MRI process:
\begin{itemize}
    \item \textbf{MRI safety questionnaire:} A detailed, mandatory pre-scan checklist to identify metallic foreign bodies, implants, or devices. Magnetic items can be dislodged or heated by the scanner, causing severe injury.
    \item \textbf{eGFR blood test:} Required before administering gadolinium contrast to patients with known kidney disease or risk factors, as a low eGFR (under 30 mL/min) increases the risk of NSF.
    \item \textbf{Orbital X-ray:} Recommended for patients with a history of metalwork or welding who may have occult metallic foreign bodies in their eyes.
    \item \textbf{Implant card verification:} Reviewing manufacturer documentation for pacemakers or implants to confirm they are "MRI Conditional" and setting appropriate scanner parameters.
    \item \textbf{Anxiety and sedation prep:} Arranging oral or intravenous sedation for patients with severe claustrophobia.
\end{itemize}

\section*{Management}
The execution of an MRI involves patient positioning, monitoring, and specific imaging sequences:
\begin{itemize}
    \item \textbf{Surgical and device planning:} Prior to scanning, pacemakers must be programmed to "MRI mode" by a cardiac technician. Any external metallic objects (jewellery, clothing with zippers, hearing aids) must be removed.
    \item \textbf{Scanning process:} The patient lies on a motorised bed inside the cylindrical bore. Earplugs or headphones are provided to block the loud knocking noises caused by rapid switching of gradient coils. The patient must remain completely still to prevent motion artifacts.
    \item \textbf{Contrast administration:} If indicated, gadolinium is injected through an intravenous cannula during specific phases of the scan.
    \item \textbf{Post-scan management:} Patients who received sedation must be monitored until recovered and are advised not to drive. Contrast-enhanced patients are monitored for 15--30 minutes for delayed allergic reactions.
\end{itemize}

\section*{Complications}
Potential complications associated with MRI scans include:
\begin{itemize}
    \item \textbf{Projectile injury (missile effect):} Ferromagnetic objects (scissors, oxygen cylinders) brought into the scanner room being pulled violently toward the magnet, causing severe trauma.
    \item \textbf{Thermal burns:} RF pulses heating metallic implants, leads, or skin tattoos containing iron-oxide pigments, causing skin burns.
    \item \textbf{Nephrogenic Systemic Fibrosis (NSF):} A rare but fatal systemic fibrotic condition of the skin, muscles, and internal organs in patients with severe renal failure exposed to gadolinium.
    \item \textbf{Implant malfunction:} Deactivation or resetting of pacemakers, implantable cardioverter-defibrillators (ICDs), or insulin pumps due to the strong magnetic field.
    \item \textbf{Contrast extravasation:} Leakage of gadolinium into surrounding tissues at the injection site, causing localised swelling and pain.
    \item \textbf{Allergic reactions:} Rare, mild-to-moderate urticaria or extremely rare anaphylactic shock in response to gadolinium.
\end{itemize}

\section*{Prognosis and Outcomes}
The diagnostic utility of MRI is exceptional, providing unmatched detail that frequently leads to accurate diagnoses and optimised treatment plans. For example, MRI can confirm multiple sclerosis plaques, identify the exact location of a spinal disc herniation, or determine the stage of a pelvic tumour. Visualising soft-tissue injuries (such as meniscus or rotator cuff tears) allows for precise orthopaedic surgical planning. The scan itself carries no long-term biological risks, and patients who do not receive sedation can immediately resume normal daily activities.

\section*{Prevention and Self-Care}
While patients cannot prevent the need for diagnostic imaging, they can ensure a safe and comfortable MRI experience:
\begin{itemize}
    \item \textbf{Complete the safety questionnaire honestly:} Disclose all surgeries, implants, or injuries involving metal, even if they occurred decades ago.
    \item \textbf{Wear comfortable, metal-free clothing:} Wear loose clothing without metal buttons, zippers, underwires, or metallic fibres (which are sometimes found in athletic wear).
    \item \textbf{Practice breathing and relaxation:} Practice taking slow, deep breaths to help stay calm and still inside the scanner.
    \item \textbf{Use hearing protection:} Always wear the provided earplugs or headphones to protect your ears from the loud acoustic noise.
    \item \textbf{Inform staff of kidney disease:} Notify the clinic if you have chronic kidney disease, are on dialysis, or have recently received contrast.
\end{itemize}

\section*{Sources and Further Reading}
The following reputable organisations provide further information on MRI scans:
\begin{itemize}
    \item InsideRadiology (RANZCR). \textit{MRI patient information and safety guidelines.} https://www.insideradiology.com.au/
    \item Australian Department of Health and Aged Care. \textit{Medicare eligibility for MRI scans.} https://www.health.gov.au/
    \item Better Health Channel. \textit{MRI scan: what happens and how to prepare.} https://www.betterhealth.vic.gov.au/
    \item Mayo Clinic. \textit{MRI (magnetic resonance imaging) overview.} https://www.mayoclinic.org/
    \item NHS. \textit{MRI scan: uses, procedure, and safety.} https://www.nhs.uk/
\end{itemize>
